% ============================================================
% DEMO 1 — Minimal LaTeX document
% Live demo: build this from scratch, step by step
%
% Compile with:  xelatex paper.tex
% ============================================================

\documentclass[12pt, a4paper]{article}

% ----- Packages -----
\usepackage{amsmath, amssymb}   % math
\usepackage{graphicx}           % figures
\usepackage{booktabs}           % publication tables
\usepackage{hyperref}           % clickable links

% ----- Title block -----
\title{A Minimal LaTeX Document}
\author{Imad El Badisy}
\date{\today}

% ============================================================
\begin{document}

\maketitle

% ----- Abstract -----
\begin{abstract}
  This is a minimal LaTeX document built live during the demo.
  We will add sections, math, a figure, and a table step by step.
\end{abstract}

% ----- Introduction -----
\section{Introduction}

LaTeX is a document preparation system that separates
\textbf{content} from \textbf{formatting}.
You write plain text with markup commands, and LaTeX
typesets a professional PDF.

% TODO (live): add a sentence with \emph{}, \texttt{}, inline math

% ----- Methods -----
\section{Methods}

\subsection{Statistical Model}

Let $Y_i$ be the outcome for subject $i = 1, \ldots, n$.
We assume the linear model:

\begin{equation}
  Y_i = \beta_0 + \beta_1 X_i + \varepsilon_i,
  \quad \varepsilon_i \sim N(0, \sigma^2)
  \label{eq:lm}
\end{equation}

The ordinary least squares estimator is:

\begin{equation}
  \hat{\boldsymbol{\beta}} =
  (\mathbf{X}^\top \mathbf{X})^{-1} \mathbf{X}^\top \mathbf{y}
  \label{eq:ols}
\end{equation}

As shown in Equation~\ref{eq:ols}, $\hat{\boldsymbol{\beta}}$
is a linear function of $\mathbf{y}$.

% TODO (live): add aligned equations with \begin{align}

\subsection{Lists}

The data preprocessing steps are:

\begin{enumerate}
  \item Remove duplicate records
  \item Impute missing values by median
  \item Standardize continuous predictors
\end{enumerate}

% ----- Results -----
\section{Results}

\subsection{Descriptive Statistics}

Table~\ref{tab:desc} shows the baseline characteristics
of the study population.

\begin{table}[htbp]
  \centering
  \caption{Baseline characteristics of study participants.}
  \label{tab:desc}
  \begin{tabular}{lrr}
    \toprule
    Variable      & Mean  & SD    \\
    \midrule
    Age (years)   & 45.2  & 12.1  \\
    BMI           & 27.3  &  4.8  \\
    Systolic BP   & 128.4 & 14.2  \\
    \midrule
    \multicolumn{3}{l}{\small $n = 250$ participants} \\
    \bottomrule
  \end{tabular}
\end{table}

% ----- Discussion -----
\section{Discussion}

We have demonstrated how to write a minimal but complete
scientific document in \LaTeX.

% ----- Conclusion -----
\section{Conclusion}

\LaTeX{} provides precise, reproducible typesetting
for scientific manuscripts.

\end{document}
